% Historical source (removed by author request): references/_OLD_Übungsaufgaben/Übungsblatt 2 - Lineare Gleichungssysteme/Blatt 2.tex:285
\begin{ExerciseBlock}
\ExerciseHeading{09}{Textaufgaben}{linear_systems_word_problems}
Überführen Sie die folgenden Textaufgaben in lineare Gleichungssysteme und lösen Sie diese:
\begin{enumerate}[label=\alph*)]
	
	\item Die Summe zweier Zahlen beträgt 35. Ihre Differenz ist 17. Wie lauten die beiden Zahlen? 
	
	\item Sie kaufen einige Produkte ein, von denen ein Teil mit einer MwSt. von 19 $\%$ und andere mit dem ermä\ss igten Satz von  7$\%$ belastet sind. Die Nettosumme (ohne MwSt.) Ihres Einkaufs beträgt 200€, die Bruttosumme (mit MwSt.) ist 220€. Wie groß ist die Nettosumme aller Produkte mit ermä\ss igtem Mehrwertsteuersatz?
	
	\item Für Düngeversuche soll aus den 3 Düngemittelsorten A, B und C, Blumendünger gemischt werden. Die Zusammensetzung der 3 Düngemittelsorten ergibt sich aus der folgenden Tabelle:
	
	\begin{center}
	\begin{tabular}{c c c c} 
	 & A & B & C \\ 
	\hline 
	Kalium     & $40\%$ & $30\%$ & $50\%$ \\
	Stickstoff & $50\%$ & $20\%$ & $30\%$ \\
	Phosphor   & $10\%$ & $50\%$ & $20\%$
	\end{tabular}
	\end{center}
	Berechnen Sie das benötigte Mischungsverhältnis der Düngemittelsorten, um ein Düngemittel mit der Zusammensetzung $40\%$ Kalium, $35\%$ Stickstoff und $25\%$ Phosphor zu erhalten. \\
	Hinweise: Erstellen Sie ein LGS, bei dem eine Gleichung jeweils einen der Inhaltsstoffe, d.h. Kalium, Stickstoff oder Phosphor, bilanziert.
	
\end{enumerate}

\end{ExerciseBlock}
\UseExerciseSolution{linear_systems_word_problems}
